%------------------------------------------------------------------------------%
\documentclass[fleqn,utf8,aspectratio=169,12pt]{beamer}

\usepackage{gaal}

\title{Posto de Uma Matriz}

%------------------------------------------------------------------------------%
\begin{document}

\addtitle

%------------------------------------------------------------------------------%
\section{Posto de uma matriz}

%------------------------------------------------------------------------------%
\begin{frame}{Posto de Uma Matriz}
  \emph{Posto} de uma matriz
  \pause
  \begin{itemize}
    \item o \emph{número de pivôs} que aparecem em sua forma escalonada
    \pause
    \item o \emph{número de linhas não nulas} de uma matriz em forma escalonada
  \end{itemize}

  \pause
  \medskip
  \[
    \posto(A)
    \pause
    \;=\;
    \rank(A)
    \pause
    \;=\;
    \text{número de pivôs de }A
  \]

  \pause
  \medskip
  Veremos adiante que o posto é o número de linhas independentes da matriz
\end{frame}

%------------------------------------------------------------------------------%
%------------------------------------------------------------------------------%
\begin{question}
Determine o posto de $A$
\[
  A =
  \begin{bmatrix}
    1 & 2 & -1 & 3  \\
    0 & 1 & 4  & 2  \\
    0 & 0 & 2  & -1 \\
    0 & 0 & 0  & 0
  \end{bmatrix}
\]
\end{question}

%------------------------------------------------------------------------------%
\begin{resolution}{Solução}
\[
  A =
  \begin{bmatrix}
    \emph{1} & 2        & -1       & 3  \\
    0        & \emph{1} & 4        & 2  \\
    0        & 0        & \emph{2} & -1 \\
    0        & 0        & 0        & 0
  \end{bmatrix}
\]
\pause
Pivôs
\[
  \pause
  a_{11} = 1
  \pause  \qquad\qquad
  a_{22} = 1
  \pause  \qquad\qquad
  a_{33} = 2
\]
\pause
Três pivôs
\[
  \pause
  \posto(A) = 3
\]
\end{resolution}

%------------------------------------------------------------------------------%
\endinput

%------------------------------------------------------------------------------%
\begin{question}
  Determine o posto de $B$
  \[
    B =
    \begin{pmatrix}
      1 & 2 & 0 & -1 & 3 \\
      0 & 0 & 1 & 2  & 4 \\
      0 & 0 & 0 & 0  & 0
    \end{pmatrix}
  \]
\end{question}

%------------------------------------------------------------------------------%
\begin{resolution}{Solução}
  \[
    B =
    \begin{pmatrix}
      1 & 2 & 0 & -1 & 3 \\
      0 & 0 & 1 & 2  & 4 \\
      0 & 0 & 0 & 0  & 0
    \end{pmatrix}
  \]

  \pause
  A matriz está na forma escalonada
  \pause
  e possui duas linhas não nulas
  \[
    \pause
    \posto(B) = 2
  \]

\end{resolution}

%------------------------------------------------------------------------------%
\endinput


%------------------------------------------------------------------------------%
\begin{frame}{Posto e Variáveis Livres}
  Dado o sistema linear com \emph{$n$ incógnitas}
  \[
    A_{m\times n} \; \mathbf{x}_{n\times 1} = \mathbf{b}_{m\times 1}
  \]
  \pause
  \(
    \emph{\posto(A) = r}
  \)\;
  \pause
  existem $r$ variáveis associadas a pivôs

  \medskip
  \pause
  As demais
  \(
    \emph{n - r}
  \)
  variáveis são livres,
  \pause
  desde que o sistema seja compatível
\end{frame}

%------------------------------------------------------------------------------%
%------------------------------------------------------------------------------%
\begin{question}
  Determine o número de variáveis livres do sistema linear associado à matriz
  \[
    A =
    \begin{pmatrix}
      1 & 2 & 0 & -1 & 3 \\
      0 & 0 & 1 & 4  & 2 \\
      0 & 0 & 0 & 0  & 0
    \end{pmatrix}
  \]
\end{question}

%------------------------------------------------------------------------------%
\begin{resolution}{Solução}
  \[
    A =
    \begin{pmatrix}
      \emph{1} & 2 & 0        & -1 & 3 \\
      0        & 0 & \emph{1} & 4  & 2 \\
      0        & 0 & 0        & 0  & 0
    \end{pmatrix}
  \]
  \pause
  \[
    \posto(A) = 2
  \]

  \pause
  $A$ possui 5 colunas,
  \pause
  o sistema linear tem \emph{$n = 5$} incógnitas
  \[
    \pause
    \text{Variáveis livres}
    \pause
    \;=\;
    \posto(A) - n
    \pause
    \;=\;
    5 - 2
    \pause
    \;=\;
    3
  \]
\end{resolution}

%------------------------------------------------------------------------------%
\endinput


%------------------------------------------------------------------------------%
\begin{frame}{Posto de Uma Matriz Quadrada}
  Para uma matriz quadrada
  \(
    A\in\mathbb{R}^{n\times n}
  \)

  \pause
  temos
  \[
    \posto(A) \leq n
  \]

  \pause
  \bigskip
  Quando
  \;\(
    \posto(A) = n
  \)\;
  \pause
  dizemos que $A$ possui \emph{posto completo}
\end{frame}

%------------------------------------------------------------------------------%
%------------------------------------------------------------------------------%
\begin{question}
  Analise o posto da matriz
  \quad
  \(
    A =
    \begin{pmatrix}
      1 & 2 & 0 \\
      0 & 1 & 3 \\
      0 & 0 & 2
    \end{pmatrix}
  \)
\end{question}

%------------------------------------------------------------------------------%
\begin{resolution}{Solução}
  \[
    A=
    \begin{pmatrix}
      \emph{1} & 2        & 0        \\
      0        & \emph{1} & 3        \\
      0        & 0        & \emph{2}
    \end{pmatrix}
  \]
  \pause
  Três pivôs
  \[
    \pause
    \posto(A) = 3
  \]
  \pause
  $A$ é uma matriz $3 \times 3$

  \pause
  $A$ é quadrada com $n=3$
  \pause
  e \(\posto(A) = 3\)

  \pause
  $A$ possui \emph{posto completo}
\end{resolution}

%------------------------------------------------------------------------------%
\endinput

%------------------------------------------------------------------------------%
\begin{question}
  Analise o posto da matriz
  \quad
  \(
    B =
    \begin{pmatrix}
      1 & 2 & 3 \\
      0 & 1 & 4 \\
      0 & 0 & 0
    \end{pmatrix}
  \)
\end{question}

%------------------------------------------------------------------------------%
\begin{resolution}{Solução}
  \[
    B =
    \begin{pmatrix}
      \emph{1} & 2        & 3 \\
      0        & \emph{1} & 4 \\
      0        & 0        & 0
    \end{pmatrix}
  \]

  \pause
  \[
    \posto(B) = 2
  \]
  \pause
  $B$ é uma matriz $3 \times 3$
  \[
    \pause
    \posto(B)
    \;=\;
    2
    \pause
    \;<\;
    3
    \;=\;
    n
  \]
  \pause
  $B$ não possui posto completo
\end{resolution}

%------------------------------------------------------------------------------%
\endinput


%------------------------------------------------------------------------------%
\begin{frame}{Teorema Fundamental para Sistemas Homogêneos}
  Dado um sistema homogêneo com $n$ incógnitas
  \[
    A \mathbf{x} = \mnula
  \]

  \pause
  \medskip
  Se
  \;\(
    \posto(A) = n
  \)\;
  \pause
  o sistema possui somente a solução trivial

  \pause
  \medskip
  Se
  \:\(
    \posto(A) < n
  \)\:
  \pause
  existem variáveis livres e   soluções não triviais
\end{frame}

%------------------------------------------------------------------------------%
%\section{Lista Mínima}
%
%%------------------------------------------------------------------------------%
%\begin{frame}{Lista Mínima}
%  \begin{enumerate}
%    \item
%  \end{enumerate}
%
%  \vfill
%  Atenção: A prova é baseada no livro, não nas apresentações
%\end{frame}

%------------------------------------------------------------------------------%
\end{document}
%------------------------------------------------------------------------------%
